% ── Section III: Methodology ─────────────────────────────────────────
% Length  : ~2 pages
% Must include:
%   - Fig. 1: pipeline overview (reference before describing)
%   - C1: equations for TrafficState, parameter mapping
%   - C2: CTRV state vector + predict equations
%   - C3: reliability score formula + dual-threshold logic
%
% Writing tip: describe WHAT each component does before HOW.
%              Reviewers must understand the design decision first.
% ─────────────────────────────────────────────────────────────────────

\section{Methodology}
\label{sec:method}

% ── 3.0 Overview ────────────────────────────────────────────────────
\subsection{System Overview}
\label{subsec:overview}

Fig.~\ref{fig:pipeline} illustrates the proposed pipeline.
% TODO: describe the data flow in 3–4 sentences
%       Camera → Context Estimator → Adaptive Controller →
%       YOLOv8 → Modified ByteTrack → Track Reliability Filter →
%       [Accident Detection | MAPPO Observations]

\begin{figure}[t]
  \centering
  % TODO: replace with actual figure
  \includegraphics[width=\linewidth]{figures/fig1_pipeline_overview.pdf}
  \caption{Overview of the proposed context-aware adaptive perception
           pipeline. Three novel modules (C1–C3) are highlighted in
           blue.}
  \label{fig:pipeline}
\end{figure}


% ── 3.1 C1: Context Estimator ───────────────────────────────────────
\subsection{C1: Context Estimator and Adaptive Parameter Controller}
\label{subsec:c1}

% [Design rationale]
% TODO: 2–3 sentences explaining why fixed thresholds fail

% [Context Estimator]
% TODO: define TrafficState formally
% Suggested notation:
%   \mathcal{S}_t = \{\rho_t,\, h_t,\, \beta_t,\, \phi_t\}
%   where:
%     \rho_t  = vehicle density  (number of active tracks)
%     h_t     = hour of day      (0–23)
%     \beta_t = mean frame brightness (0–255)
%     \phi_t  = recent FP rate   (tracks with age < 3 / total tracks)

% [Adaptive Parameter Controller]
% TODO: define the mapping \mathcal{S}_t → \Theta_t
%   \Theta_t = \{c_{\text{conf}},\, L_{\text{buf}},\, \mathcal{M}\}
%   where:
%     c_conf  = YOLO confidence threshold
%     L_buf   = ByteTrack lost-track buffer length
%     M       = motion model selector ∈ {ctrv, standard, conservative}
%
%   Table of states → parameters (dense_peak, night, default, etc.)


% ── 3.2 C2: Class-Conditioned Motion Model ──────────────────────────
\subsection{C2: Class-Conditioned CTRV Motion Model}
\label{subsec:c2}

% [Design rationale]
% TODO: explain why linear Kalman fails for motorcycles (weaving, omega)

% [CTRV state vector]
% \mathbf{x} = [x,\; y,\; v,\; \theta,\; \omega]^\top

% [Prediction equations — use align environment]
% \begin{align}
%   x_{k+1}      &= x_k + v_k \cos\theta_k \cdot \Delta t  \label{eq:ctrv_x} \\
%   y_{k+1}      &= y_k + v_k \sin\theta_k \cdot \Delta t  \label{eq:ctrv_y} \\
%   v_{k+1}      &= v_k                                     \label{eq:ctrv_v} \\
%   \theta_{k+1} &= \theta_k + \omega_k \cdot \Delta t     \label{eq:ctrv_theta} \\
%   \omega_{k+1} &= \omega_k                                \label{eq:ctrv_omega}
% \end{align}

% [Integration into ByteTrack]
% TODO: explain that motorcycle-class tracks use CTRV;
%       car/truck tracks keep original linear Kalman.
%       Reference byte_tracker.py integration point.

% [Algorithm pseudocode — optional but recommended]
% \begin{algorithm}
%   \caption{Class-Conditioned Motion Model Selection}
%   \label{alg:motion_model}
%   \begin{algorithmic}[1]
%     \For{each active track $\tau$}
%       \If{$\tau.\text{class} = \texttt{motorcycle}$}
%         \State apply CTRV prediction (Eq.~\ref{eq:ctrv_x}–\ref{eq:ctrv_omega})
%       \Else
%         \State apply standard linear Kalman prediction
%       \EndIf
%     \EndFor
%   \end{algorithmic}
% \end{algorithm}


% ── 3.3 C3: Dual-Task Track Reliability Score ───────────────────────
\subsection{C3: Dual-Task Track Reliability Score}
\label{subsec:c3}

% [Design rationale]
% TODO: explain the precision–recall trade-off for two tasks:
%   - Accident detection → high precision → filter low-score tracks
%   - MAPPO observations → high recall   → keep medium-score tracks

% [Score formula]
% \begin{equation}
%   S(\tau) = w_1 \cdot \text{age}(\tau)
%           + w_2 \cdot \text{consistency}(\tau)
%           + w_3 \cdot \text{smoothness}(\tau)
%           + w_4 \cdot \text{class\_stability}(\tau)
%   \label{eq:reliability}
% \end{equation}
% where w_1+w_2+w_3+w_4 = 1  (TODO: justify weight values empirically)

% [Component definitions]
% TODO: define each term:
%   age(\tau)             = number of consecutive frames with confirmed detection
%   consistency(\tau)     = detection hit rate over last N frames
%   smoothness(\tau)      = inverse of IoU variance in predicted vs detected boxes
%   class_stability(\tau) = fraction of frames with consistent class label

% [Threshold split]
% TODO: define \tau_acc and \tau_mappo
%   Accident branch: S(\tau) > \tau_acc   (high threshold)
%   MAPPO branch:    S(\tau) > \tau_mappo (lower threshold)
%   Ablation will determine optimal \tau_acc and \tau_mappo
